% Work/MAT -- přepracovaná hodina 16: Operace s intervaly.
% Upravený plán bloku: 12 Množina, 13 Podmnožiny a operace,
% 14 Vennovy diagramy, 15 Vlastnosti operací, 16 Operace s intervaly,
% 17 Slovní úlohy. Původní rozpis má toto učivo v hodině 18;
% před definitivním nasazením je třeba upravit ročníkový rozpis.
% Hlavní zdroje: ucebnice/01_Zakladni_poznatky/
% Ciselne_obory/Číselné obory.tex (Intervaly);
% Mnoziny/Množiny.tex (Množinové operace).
% Návaznost na 1_rocnik/pripravy/08_CIS_Intervaly.tex a
% 18_MNZ_OperaceSIntervaly.tex (původní lekce).
% Metodická inspirace: Realisticky.cz, 010308 Intervaly,
% zejména hraniční případy dotýkajících se intervalů a práce na ose;
% https://www.realisticky.cz/ucebnice/01%20Matematika%20S%C5%A0/
% 01%20Z%C3%A1kladn%C3%AD%20poznatky/03%20Mno%C5%BEiny/08%20Intervaly.pdf
% Příklady a grafy nově sestavené; nejsou doslovnými přepisy sbírky.
\documentclass[11pt,a4paper]{article}
% Spolecny styl pro ucitelske pripravy MAT (Work/MAT)
% Verze 1.1, 2026-10-23
\NeedsTeXFormat{LaTeX2e}
\ProvidesPackage{priprava}[2026/10/23 v1.1 Sablona ucitelskych priprav MAT]

% Zakladni typografie a cestina
\RequirePackage[T1]{fontenc}
\RequirePackage[utf8]{inputenc}
\RequirePackage[czech]{babel}
\RequirePackage{lmodern}
\RequirePackage{microtype}
\RequirePackage{amsmath,amssymb}
\RequirePackage{xcolor}
\RequirePackage{enumitem}
\RequirePackage[most]{tcolorbox}
\RequirePackage{tabularx,booktabs,array}
\RequirePackage{fancyhdr}
\RequirePackage{needspace}
\RequirePackage{geometry}
\RequirePackage[hidelinks]{hyperref}

\geometry{a4paper,top=17mm,bottom=18mm,left=23mm,right=23mm}
\definecolor{notegray}{gray}{0.39}
\setlength{\parindent}{0pt}
\setlength{\parskip}{2.5pt plus 1pt minus 0.5pt}
\setlist[itemize]{topsep=2pt,itemsep=1pt,parsep=0pt,leftmargin=1.45em}
\setlist[enumerate]{topsep=2pt,itemsep=2pt,parsep=0pt,leftmargin=1.65em}
\setcounter{secnumdepth}{0}
\renewcommand{\arraystretch}{1.12}

% Zahlavi: vlevo rocnik | globalni tema, vpravo cislo: nazev hodiny
% Zapati zustava zcela prazdne.
\pagestyle{fancy}
\fancyhf{}
\renewcommand{\headrulewidth}{0.35pt}
\renewcommand{\footrulewidth}{0pt}
\setlength{\headheight}{14pt}
\newcommand{\nastavPripravu}[4]{%
  \fancyhead[L]{\small #1\enspace\textbar\enspace #2}%
  \fancyhead[R]{\small #3: #4}%
}

% Minimalisticke cernobile boxy. Kazda definice zvlast.
\tcbset{commonlesson/.style={enhanced,breakable=false,colback=white,colframe=black,
  colbacktitle=white,coltitle=black,boxrule=0.45pt,
  arc=0.8mm,outer arc=0.8mm,left=3mm,right=3mm,top=1.6mm,bottom=1.6mm,
  toptitle=1mm,bottomtitle=0.7mm,fonttitle=\bfseries,
  attach boxed title to top left={xshift=2.5mm,yshift=-1.15mm},
  boxed title style={colback=white,colframe=white,boxrule=0pt,left=1mm,right=1mm}}}
\newtcolorbox{zadaniBox}[1][]{commonlesson,title={Zadání},#1}
\newtcolorbox{prikladBox}[1][]{commonlesson,title={Příklad},#1}
\newtcolorbox{definiceBox}[1][]{commonlesson,title={Definice},#1}

% Spolecny box pro zadani a jeho reseni, bez mezery mezi castmi.
% Pouziti:
%   \begin{prikladReseniBox}
%     Zadani ulohy.
%     \tcblower
%     \begin{align*} ... \end{align*}
%   \end{prikladReseniBox}
% Spodni cast ma jemny sedy podklad a delici linku; lze v ni sazet
% vice radku matematiky. Obe casti maji jen jeden spolecny ramecek.
\tcbset{solutionlesson/.style={
  commonlesson,
  enhanced jigsaw,
  breakable,
  colbacklower=black!3,
  segmentation style={solid,draw=black!25,line width=0.35pt},
  fontlower=\small\color{notegray},
  before lower={\noindent\textbf{Řešení:}\enspace},
  leftlower=3mm,rightlower=3mm
}}
\newtcolorbox{prikladReseniBox}[1][]{solutionlesson,title={Příklad},#1}
\newtcolorbox{zadaniReseniBox}[1][]{solutionlesson,title={Zadání},#1}

% Pomucky pro autora pripravy: sedive poznamky; \reseni pro starsi pripravy
\newcommand{\poznamka}[1]{\par\begingroup\small\color{notegray}\itshape
  \textnormal{\textbf{Poznámka pro učitele.}}\enspace #1\par\endgroup}
\newcommand{\reseni}[1]{\par\begingroup\small\color{notegray}
  \textbf{Očekávané řešení:}\enspace #1\par\endgroup}
\newcommand{\blok}[1]{\needspace{5\baselineskip}\subsection*{#1}}
\newcommand{\vyrok}[1]{\uv{#1}}

\usepackage{tikz}
\usetikzlibrary{arrows.meta,calc}
\nastavPripravu{1. ročník}{Množiny}{16}{Operace s intervaly}

% Pomocné kreslicí prvky černobílých číselných os.
\newcommand{\plny}[2]{\fill[black] (#1,#2) circle (2.65pt);}
\newcommand{\prazdny}[2]{\draw[thick,fill=white] (#1,#2) circle (2.65pt);}
\newcommand{\pruh}[3]{\draw[line width=1.55pt] (#1,#3)--(#2,#3);}
\newcommand{\popisek}[2]{\node[anchor=east] at (#1,#2) {\small};}
\newcommand{\osacisla}[2]{%
  \draw[-{Latex[length=2mm]},thin] (#1-.36,0)--(#2+.45,0) node[right] {$x$};%
  \foreach \t in {#1,...,#2}{\draw (\t,.075)--(\t,-.075) node[below=3pt,font=\scriptsize] {$\t$};}%
}
\begin{document}

\blok{Zpátky na číselnou osu: co mají dvě podmínky společného?}
Na 8. hodině jsme zavedli intervaly, na 13. hodině operace s množinami.
Nyní obojí propojíme: \textbf{interval je množina reálných čísel}.
Průnik proto znamená \uv{splň obě podmínky}, sjednocení \uv{splň alespoň jednu},
rozdíl \uv{splň první, ale ne druhou}.

\begin{definiceBox}
Pro intervaly $A,B\subseteq\mathbb R$ platí stejné množinové definice:
\[
\begin{aligned}
A\cap B&=\{x\in\mathbb R\mid x\in A\ \land\ x\in B\},\\[-1mm]
A\cup B&=\{x\in\mathbb R\mid x\in A\ \lor\ x\in B\},\\[-1mm]
A\setminus B&=\{x\in\mathbb R\mid x\in A\ \land\ x\notin B\}.
\end{aligned}
\]
\end{definiceBox}
\poznamka{Definice jsou připomenutím, nikoli novým výkladem z 13. hodiny.
Na tabuli nejprve nech žáky odhadnout, zda výsledek bude jeden interval,
jediný bod, prázdná množina, nebo sjednocení více intervalů.}

\blok{Průnik a sjednocení překrývajících se intervalů}
\begin{prikladReseniBox}[title={P1: jedna osa, čtyři řádky}]
Nechť $A=\langle-3;2)$ a $B=(0;5\rangle$. Urči $A\cap B$ a $A\cup B$.
Vyznač oba vstupní intervaly i výsledky na číselné ose.
\tcblower
\begin{center}
\begin{tikzpicture}[x=.91cm,y=.66cm]
  \osacisla{-4}{6}
  \node[left] at (-4.45,3.65) {$A$}; \pruh{-3}{2}{3.65} \plny{-3}{3.65} \prazdny{2}{3.65}
  \node[left] at (-4.45,2.60) {$B$}; \pruh{0}{5}{2.60} \prazdny{0}{2.60} \plny{5}{2.60}
  \node[left] at (-4.45,1.55) {$A\cap B$}; \pruh{0}{2}{1.55} \prazdny{0}{1.55} \prazdny{2}{1.55}
  \node[left] at (-4.45,.50) {$A\cup B$}; \pruh{-3}{5}{.50} \plny{-3}{.50} \plny{5}{.50}
\end{tikzpicture}
\end{center}
\[
\boxed{A\cap B=(0;2)},\qquad
\boxed{A\cup B=\langle-3;5\rangle}.
\]
U průniku musí číslo ležet v \emph{obou} intervalech:
$0$ je vyloučené z $B$ a $2$ z $A$.
U sjednocení mezi $-3$ a $5$ žádný bod nechybí;
oba vnější kraje patří alespoň do jedné množiny.
\end{prikladReseniBox}

\poznamka{Nech žáky před zobrazením řešení jednotlivě otestovat
$-3$, $0$, $2$ a $5$. Pouhé porovnání čtyř krajních čísel nestačí:
rozhodují i otevřené či uzavřené konce.}

\begin{zadaniReseniBox}[title={Volitelná diagnostika P1: příslušnost konkrétních čísel},breakable=false]
Předcházející výsledek ověř pro čísla $-3,0,1,2,5$.
Rozhodni, zda leží v $A$, $B$, v jejich průniku a ve sjednocení.
\tcblower
\begin{center}
\begin{tabular}{c|ccccc}
\toprule
$x$ & $-3$ & $0$ & $1$ & $2$ & $5$\\\midrule
$x\in A$ & ano & ano & ano & ne & ne\\
$x\in B$ & ne & ne & ano & ano & ano\\
$x\in A\cap B$ & ne & ne & ano & ne & ne\\
$x\in A\cup B$ & ano & ano & ano & ano & ano\\
\bottomrule
\end{tabular}
\end{center}
Průnik vyžaduje dvě odpovědi \uv{ano}, sjednocení alespoň jednu.
\end{zadaniReseniBox}

\blok{Tři typy dotyku: bod, díra, navazující interval}
\begin{prikladReseniBox}[title={P2: jediná změna závorky změní výsledek},breakable=false]
Porovnej následující tři dvojice. V každé urči průnik a sjednocení:
\begin{enumerate}[label=\alph*)]
\item $C=(-2;1\rangle$, $D=\langle1;4)$;
\item $E=(-2;1)$, $F=(1;4)$;
\item $G=(-2;1)$, $H=\langle1;4)$.
\end{enumerate}
\tcblower
\begin{center}
\begin{tikzpicture}[x=.94cm,y=.62cm]
\osacisla{-3}{5}
\node[left] at (-3.5,3.9) {a)};
\pruh{-2}{1}{3.9} \prazdny{-2}{3.9} \plny{1}{3.9}
\pruh{1}{4}{3.35} \plny{1}{3.35} \prazdny{4}{3.35}
\node[left] at (-3.5,2.55) {b)};
\pruh{-2}{1}{2.55} \prazdny{-2}{2.55} \prazdny{1}{2.55}
\pruh{1}{4}{2.00} \prazdny{1}{2.00} \prazdny{4}{2.00}
\node[left] at (-3.5,1.22) {c)};
\pruh{-2}{1}{1.22} \prazdny{-2}{1.22} \prazdny{1}{1.22}
\pruh{1}{4}{.67} \plny{1}{.67} \prazdny{4}{.67}
\end{tikzpicture}
\end{center}
\begin{align*}
\text{a)}\quad &C\cap D=\{1\}, & C\cup D&=(-2;4),\\
\text{b)}\quad &E\cap F=\varnothing, & E\cup F&=(-2;1)\cup(1;4),\\
\text{c)}\quad &G\cap H=\varnothing, & G\cup H&=(-2;4).
\end{align*}
\textbf{Pozor:} v c) je průnik prázdný, ale sjednocení tvoří jeden interval,
protože bod $1$ obsahuje $H$. V b) bod $1$ chybí v obou intervalech.
V a) je průnik jednoprvková množina $\{1\}$, nikoli samotné číslo $1$.
\end{prikladReseniBox}

\blok{Rozdíl: nezapomeň, že záleží na pořadí}
\begin{prikladReseniBox}[title={P3: odebrání prostřední části}]
Pro $A=\langle-2;4)$ a $B=(1;6\rangle$ vypočti oba rozdíly.
\tcblower
\begin{center}
\begin{tikzpicture}[x=.91cm,y=.64cm]
\osacisla{-3}{7}
\node[left] at (-3.48,3.55) {$A$}; \pruh{-2}{4}{3.55} \plny{-2}{3.55} \prazdny{4}{3.55}
\node[left] at (-3.48,2.60) {$B$}; \pruh{1}{6}{2.60} \prazdny{1}{2.60} \plny{6}{2.60}
\node[left] at (-3.48,1.65) {$A\setminus B$}; \pruh{-2}{1}{1.65} \plny{-2}{1.65} \plny{1}{1.65}
\node[left] at (-3.48,.70) {$B\setminus A$}; \pruh{4}{6}{.70} \plny{4}{.70} \plny{6}{.70}
\end{tikzpicture}
\end{center}
\[
\boxed{A\setminus B=\langle-2;1\rangle},\qquad
\boxed{B\setminus A=\langle4;6\rangle}.
\]
Jednička zůstává v prvním rozdílu, protože není v $B$.
Čtyřka patří do druhého rozdílu, protože není v $A$.
\end{prikladReseniBox}

\begin{prikladReseniBox}[title={P3b: když odebereme vnitřní interval},breakable=false]
Co zbude z $K=\langle-3;5\rangle$, odebereme-li $L=(0;2)$?
\tcblower
\[
\boxed{K\setminus L=\langle-3;0\rangle\cup\langle2;5\rangle}.
\]
Odebrali jsme všechna čísla \emph{mezi} $0$ a $2$.
Oba krajní body $0$ a $2$ však v $L$ nejsou, a proto zůstávají v $K$.
Výsledkem rozdílu nemusí být jediný interval.
\end{prikladReseniBox}
\poznamka{P3b je volitelné doplnění k P3. Typická chyba: označení
$0$ a $2$ za vymazané body, přestože odebíraná množina je otevřená.}

\blok{Doplněk v $\mathbb R$: co na ose zůstalo mimo interval?}
\begin{definiceBox}
\textbf{Doplněk intervalu} $I$ vzhledem k $\mathbb R$
je $I'_{\mathbb R}=\mathbb R\setminus I$.
Obecně závisí doplněk na zvolené základní množině.
\end{definiceBox}
\begin{prikladReseniBox}[title={P4: dvě neomezené části},breakable=false]
Urči doplněk intervalu $I=(-1;3\rangle$ vzhledem k $\mathbb R$.
\tcblower
\begin{center}
\begin{tikzpicture}[x=.85cm,y=.65cm]
\osacisla{-5}{6}
\node[left] at (-5.5,2.05) {$I$}; \pruh{-1}{3}{2.05} \prazdny{-1}{2.05} \plny{3}{2.05}
\node[left] at (-5.5,.75) {$I'_{\mathbb R}$};
\draw[-{Latex[length=2mm]},line width=1.55pt] (-1,.75)--(-5.05,.75);
\plny{-1}{.75}
\draw[-{Latex[length=2mm]},line width=1.55pt] (3,.75)--(5.95,.75);
\prazdny{3}{.75}
\end{tikzpicture}
\end{center}
\[
\boxed{I'_{\mathbb R}=(-\infty;-1\rangle\cup(3;+\infty)}.
\]
Číslo $-1$ do $I$ nepatří, proto je v doplňku.
Číslo $3$ v $I$ je, proto v doplňku není.
\end{prikladReseniBox}
\poznamka{V povinném jádru stačí P1, P2 (a, b, c stručně), P3 nebo P4
podle tempa a tři samostatné položky S1. Nezkoušej během 45 minut
vyčerpat celý dokument; ostatní příklady tvoří připravenou rezervu.}

\blok{Samostatná práce: od obrázku ke správným závorkám}
\begin{zadaniReseniBox}[title={S1: určuj $\cap$ a $\cup$, vždy nakresli osu}]
\begin{enumerate}[label=\alph*)]
\item $A=(-\infty;2\rangle$, $B=(-1;4)$.
\item $A=\langle0;3\rangle$, $B=(3;7\rangle$.
\item $A=(-3;0)$, $B=\langle1;4)$.
\item $A=\langle-2;5)$, $B=\langle0;3\rangle$.
\item $A=(-\infty;0)$, $B=\langle0;+\infty)$.
\end{enumerate}
\tcblower
\begin{enumerate}[label=\alph*)]
\item $A\cap B=(-1;2\rangle$, $A\cup B=(-\infty;4)$.
\item $A\cap B=\varnothing$, $A\cup B=\langle0;7\rangle$.
Bod $3$ ve sjednocení je, protože patří do $A$.
\item $A\cap B=\varnothing$, $A\cup B=(-3;0)\cup\langle1;4)$.
Mezi $0$ a $1$ je skutečná mezera.
\item $A\cap B=\langle0;3\rangle$, $A\cup B=\langle-2;5)$.
Celé $B$ leží v $A$.
\item $A\cap B=\varnothing$, $A\cup B=\mathbb R$;
$0$ obsahuje pouze $B$.
\end{enumerate}
\end{zadaniReseniBox}

\begin{zadaniReseniBox}[title={S2: rozdíl a doplněk}]
\begin{enumerate}[label=\alph*)]
\item Pro $A=\langle-3;5\rangle$, $B=(0;2)$ vypočti $A\setminus B$.
\item Urči doplněk $\langle2;6)$ v $\mathbb R$.
\item Urči $(\langle-1;4\rangle\setminus\langle1;3))\cap(-\infty;2)$.
\end{enumerate}
\tcblower
\begin{enumerate}[label=\alph*)]
\item $\langle-3;0\rangle\cup\langle2;5\rangle$:
nula ani dvojka do $B$ nepatří, takže obě zůstávají.
\item $(-\infty;2)\cup\langle6;+\infty)$.
\item Nejprve $\langle-1;1)\cup\langle3;4\rangle$, potom průnik
s $(-\infty;2)$ dává $\langle-1;1)$.
\end{enumerate}
\end{zadaniReseniBox}

\blok{Volitelná rezerva: složitější případy a kontrola oboru}
\begin{zadaniReseniBox}[title={R1: tři množiny a celá čísla}]
\begin{enumerate}[label=\alph*)]
\item Pro $A=(-\infty;3)$, $B=\langle-1;5\rangle$, $C=(1;+\infty)$
urči $(A\cap B)\cap C$ a $(A\cup B)\cap C$.
\item Urči $\langle-2;3)\cap\mathbb Z$ a $\langle-2;3)\cap\mathbb N$.
Je první výsledek reálný interval?
\item Rozhodni, zda platí
$\bigl((-\infty;2)\cup\langle2;+\infty)\bigr)=\mathbb R$.
\end{enumerate}
\tcblower
\begin{enumerate}[label=\alph*)]
\item $A\cap B=\langle-1;3)$, tedy $(A\cap B)\cap C=(1;3)$.
Dále $A\cup B=(-\infty;5\rangle$, tedy $(A\cup B)\cap C=(1;5\rangle$.
\item $\{-2,-1,0,1,2\}$ a $\{1,2\}$ při $\mathbb N=\{1,2,\ldots\}$.
Nejde o interval v $\mathbb R$, protože mezi vypsanými čísly chybí necelá čísla.
\item Ano: číslo $2$ patří do pravého intervalu, menší čísla do levého,
větší také do pravého.
\end{enumerate}
\end{zadaniReseniBox}

\begin{zadaniReseniBox}[title={R2: vzdálenost na ose -- bonus},breakable=false]
Čísla vzdálená od $2$ nejvýše o $3$ tvoří množinu $M$.
Zapiš $M$ intervalem a spočti $M\setminus(0;3)$.
\tcblower
Jde o $|x-2|\le3$, tedy $M=\langle-1;5\rangle$.
Po vyjmutí otevřeného intervalu $(0;3)$ zůstane
$\langle-1;0\rangle\cup\langle3;5\rangle$.
Bod $0$ i $3$ zůstávají, protože je nevyjímáme.
\end{zadaniReseniBox}

\blok{Domácí procvičování}
\begin{zadaniReseniBox}[title={D1: postupné ztížení, 15--20 minut}]
\begin{enumerate}[label=\alph*)]
\item Pro $A=(-4;2\rangle$, $B=\langle-1;6)$ určuj $A\cap B$, $A\cup B$, $A\setminus B$.
\item Urči průnik a sjednocení intervalů $(-\infty;-1)$ a $(-1;+\infty)$.
\item Pro $A=\langle2;4\rangle$, $B=\langle4;7\rangle$ vypočti $A\cap B$ a $A\cup B$.
\item Pro $I=(-2;3\rangle$ urči $I\cap\mathbb Z$ a $I'_{\mathbb R}$.
\item Pro $A=\langle-2;2\rangle$, $B=(0;5)$ a $C=\langle1;3\rangle$
vypočti $(A\cup B)\cap C$.
\end{enumerate}
\tcblower
\begin{enumerate}[label=\alph*)]
\item $\langle-1;2\rangle$, $(-4;6)$, $(-4;-1)$.
\item $\varnothing$; $(-\infty;-1)\cup(-1;+\infty)=\mathbb R\setminus\{-1\}$.
\item $\{4\}$; $\langle2;7\rangle$.
\item $\{-1,0,1,2,3\}$;
$(-\infty;-2\rangle\cup(3;+\infty)$.
\item $A\cup B=\langle-2;5)$, po průniku s $C$ dostaneme $\langle1;3\rangle$.
\end{enumerate}
\end{zadaniReseniBox}

\poznamka{Doporučené 45minutové jádro včetně organizace:
3 min aktivace znalostí, 9 min P1, 9 min P2, 8 min vybraný P3 nebo P4,
12 min S1 (a--c) se společnou kontrolou, 4 min shrnutí a domácí práce.
Samostatné S1 d), e), S2 i rezerva jsou volitelná nabídka.
U každého chybného řešení se vrať ke konkrétnímu krajnímu bodu a ověř jeho příslušnost.}

\blok{Volitelná příloha: prázdné osy pro tabuli a vlastní kontrolu}
\poznamka{Předlohy lze promítat nebo využít pro rychlou kontrolu řešení.
Nejde o další povinné zadání. Nech žáka zakreslit $A$, $B$ a výsledek
ve třech řádcích a vysvětlit zahrnutí každého kritického bodu.}
\begin{center}
\begin{tikzpicture}[x=.99cm,y=.7cm]
  \foreach \h/\lab in {2.4/A,1.2/B,0/{Výsledek}}{
    \node[anchor=east] at (-4.45,\h) {$\lab$};
    \draw[-{Latex[length=1.7mm]},thin] (-4.15,\h)--(6.25,\h);
    \foreach \x in {-3,-2,-1,0,1,2,3,4,5}{
      \draw (\x,\h-.07)--(\x,\h+.07);
      \node[below=3pt,font=\tiny] at (\x,\h) {$\x$};
    }
  }
\end{tikzpicture}
\end{center}

\end{document}
